\section*{Abstract}
\addcontentsline{toc}{section}{Abstract}

Read the paper "A guide to writing articles in energy science" from Martin Weiss \& Alexandra M. Newman (2011)

Abstract and Introduction
The abstract (short summary) is placed at the very beginning of the work. It comprises at most one page.
In the abstract, only the following questions are of interest: What is your work about (question)? Which methods were used (solution)? What is the outcome ( results)?
A reader who has read the abstract should know the rough outlines of the work and its most important results (no details!).
Not included in the abstract are the motivation, history of the problem, or information about the structure of the work
Usually, the abstract does not contain equations, figures, tables, literature citations, or other references to other pages of the paper
The introduction marks the beginning of the actual work. It leads to the research question by showing the bigger picture of the research and placing it in relation to existing knowledge (literature). 
In contrast to the abstract, the introduction, therefore, also includes statements on the motivation and background of the work.
The introduction then defines the scientific questions to be examined in the context of the work and mentions the methods used or further developed to solve the problem.
The last section of the introduction gives an overview of the structure of the work ("storyline"). A brief description suffices of the steps of the research described in each chapter and how the chapters build on each other. However, no results are presented yet. Only the procedure is explained!